\section{Modelado y discretización}
\subsection{Modelo identificado adoptado para el diseño}
El diseño utiliza el modelo identificado en el Laboratorio 5 mediante
mediciones entre J5 y J7:
\begin{equation}
G_p(s)=\frac{1.002941}{(0.00906915s+1)(0.000774399s+1)}.
\label{eq:modelo_identificado}
\end{equation}
El ensayo abierto de este laboratorio produjo una ganancia incremental cercana
a 0,98, coherente con la identificación previa. El modelo nominal se conserva
como referencia del circuito, mientras que el controlador y los estimadores se
calculan con el modelo de~\eqref{eq:modelo_identificado}.

Para mantener $x_1$ y $x_2$ en voltios y la salida $y=x_2$, se usa una
realización en cascada. La identificación entrada--salida no separa las
ganancias de ambas etapas; se conserva la ganancia nominal de la primera y se
asigna la segunda para reproducir exactamente la ganancia total
de~\eqref{eq:modelo_identificado}. Este reparto fija una base de estados y no
altera la transferencia medida.

\subsection{Selección de muestreo y modelo discreto}
Se adoptó $T_s=\SI{2}{\milli\second}$. Las cuatro capturas físicas registraron
una mediana de período de \SI{2000}{\micro\second}. La discretización ZOH del
modelo identificado resulta
\begin{equation}
A=\begin{bmatrix}0.80210&0\\-0.24963&0.075574\end{bmatrix},\qquad
B=\begin{bmatrix}-0.63161\\0.13046\end{bmatrix},
\end{equation}
\begin{equation}
C=\begin{bmatrix}0&1\end{bmatrix},\qquad D=0.
\end{equation}
Se verificaron
$\operatorname{rango}\mathcal C(A,B)=2$ y
$\operatorname{rango}\mathcal O(A,C)=2$. Por tanto, el modelo identificado
discreto es controlable y observable desde J7.

\subsection{Consistencia entre los tres niveles de validación}
El mismo modelo identificado y las mismas ganancias se usaron en: (i) el
cálculo lineal ideal; (ii) la simulación no ideal con cuantización ADC/DAC; y
(iii) el firmware programado para los ensayos. La referencia fue
\SI{50}{\milli\volt}, aplicada a partir de la muestra 32
($t=\SI{64}{\milli\second}$), con estados reiniciados durante el pretrigger.
Estas condiciones establecen una correspondencia directa entre cálculo,
simulación y banco.
